\section{Structure des programmes} 
%\label{sec:4:ecrir-progs} 

De même que pour écrire correctement en français il faut respecter
des règles précises d'orthographe et de grammaire, l'écriture de
programmes corrects est soumise à un certain nombre de {\em règles
de syntaxe}\index{Règle de syntaxe} heureusement plus simples que pour le français. La différence
est que des phrases françaises comportant des fautes d'orthographe ou
de grammaire peuvent rester compréhensibles par des francophones alors
qu'un programme comportant des 
{\em erreurs syntaxiques}\index{Erreur syntaxique} ne peut
pas être compilé, c'est-à-dire traduit dans le langage de l'or\-di\-na\-teur.
Cette compilation est effectuée par un programme (appellé {\em compilateur})\index{Compilateur}
prévu pour indiquer nos erreurs mais pas pour les corriger. Pour corriger
nos erreurs de manière satisfaisante, le compilateur devrait pouvoir comprendre 
nos intentions, un problème bien difficile.

Le programme {\idt Cocotte} ci-dessous est un exemple de
programme \pascal\ correct.

\begin{exemple}[Programme syntaxiquement correct]{\hfill
\label{prog:Cocotte} \input Cocotte.tex
\noindent $\Box$ }\end{exemple}

\noindent Tous les programmes \pascal\ ont la même forme :

\begin{center}
\begin{tabular}{@{}|p{12.5cm}|@{}} \hline
{\em Programme \pascal\ :}\\ \hline \hline
\begin{minipage}[t]{12.5cm}
{\medskip
\mbox{\blc\ \ \ {\mcl program}\hspace{0.3em}\blanc\hspace{0.3em}\idt NomDuProgramme$;$}\par 
\mbox{\blc\ \ \ \ \mcl uses\idt\ \ldots\ $;$}\par 
\mbox{\blc\ \ \ \mcl begin}\par 
\mbox{\blc\ \ \ \ \small\rm  \ldots\ Instructions du programme (séparées par des `$;$') \ldots}\par 
\mbox{\blc\ \ \ \mcl end\idt$.$}\par 
\medskip}
\end{minipage}\\ \hline 
\end{tabular}
\end{center}

Un programme \pascal\ doit commencer par le mot anglais {\mcl program}. 
Ensuite, on indique le {\idt  NomDuProgramme} qui est 
impérativement suivi d'un point-virgule `$;$'. 

Le {\idt  NomDuProgramme} est choisi par le programmeur
à sa convenance, de façon à évoquer ce à quoi sert son programme.
 La seule règle à res\-pec\-ter est que le {\idt  NomDuProgramme}
ne doit comporter que des lettres minuscules ou
majuscules non accentuées\footnote{\ Les accents n'existent pas en anglais
(à de très rares exceptions près, comme dans les mots `naïve', `protégé' ou
`rôle' empruntés au français).}, des chiffres ou
des soulignés `\_'\footnote{\ Le souligné `\_' ne doit pas être confondu avec le tiret `-'.}
et doit commencer par une lettre. Par exemple {\idt  S}, 
{\idt  s}, {\idt  LeStegosaure} ou {\idt  Le\_Stegosaure} sont des noms corrects de
programmes alors que {\idt  1S}, 
{\idt  s!}, {\idt  Le\ Stegosaure} ou {\idt  Le-Stégosaure} ne respectent
pas la règle d'écriture des noms de programmes. 

Il y a obligatoirement un blanc
entre le mot {\mcl program} et le {\idt  Nom\-Du\-Pro\-gram\-me}. Dans un livre il y a des blancs partout et même, en y regardant bien,
plus de blanc que toute autre chose. Mettre un blanc dans un livre pour dire qu'il faut taper
un blanc est donc idiot parce que la plupart des lecteurs ne s'intéressent pas aux 
blancs dans les livres mais simplement à ce qui y est écrit.
Nous utiliserons donc le symbole `\blanc' pour écrire nos blancs noir sur \hbox{blanc$\:$!}
Autrement dit nous écrivons \blanc\ pour dire qu'il faut taper sur la grande barre 
située en bas du clavier de l'or\-di\-na\-teur qui s'appelle la 
{\em barre d'espacement}\index{Barre d'espacement}
pour obtenir un blanc `\ ' dans le programme. 
Nous ne le faisons pas dans les programmes donnés en exemple qui sont
présentés comme ils apparaissent à l'écran. De même, à la fin d'une 
ligne il faut taper sur la touche de retour à la ligne 
\framebox[11.mm][l]{$\hookleftarrow$~~~} pour passer au début de la
ligne suivante. Nous ne l'écrivons pas dans les livres car tous les lecteurs savent
qu'il faut passer au début de la ligne suivante quand ils sont arrivés en fin de ligne. 
Certains or\-di\-na\-teurs ne font rien quand on tape une commande tant que l'on n'a pas
tapé sur la touche de retour à la ligne. Ne vous êtes-vous jamais laissé
surprendre en train d'attendre que l'or\-di\-na\-teur fasse quelque chose alors qu'ayant 
oublié de taper sur la touche de retour à la ligne, l'or\-di\-na\-teur 
attendait également que
vous \hbox{réagissiez$\:$}

Après la première ligne du programme figure une ligne commençant par
{\mcl uses} (qui signifie `utilise' en anglais). Cette ligne sert à indiquer
que nous utilisons des commandes pour le \robot\ dans le programme. Elle dépend
du type d'or\-di\-na\-teur utilisé. Son contenu exact est fourni sur la disquette
d'ac\-com\-pa\-gne\-ment du livre. 

Entre le {\mcl begin} et le {\mcl end}, on trouve les instructions 
du programme dans l'ordre où elles doivent être exécutées par
l'or\-di\-na\-teur. Ces instructions sont
séparées par des points-virgules. Les instructions du programme
peuvent être des commandes du \hbox{\robot\ :} {\idt av}, {\idt vg}, {\idt vd}, {\idt pqtd}, 
{\idt pqtg}, {\idt phtd}, {\idt phtg}, {\idt lc}, {\idt bc}, {\idt dc}, {\idt dp}, 
{\idt dg} ou {\idt st}. En général les programmes de pilotage du \robot\ se
terminent par la commande {\idt st} qui permet d'attendre que nous
tapions sur la touche de retour à la ligne avant que le dessin ne 
disparaisse à la fin de l'exécution du programme.

Finalement, le texte d'un programme \pascal\ se termine par {\mcl end}
suivi d'un point final `.'. 

Les textes entre accolades {\idt \{} et {\idt \}} sont des {\em commentaires}
ignorés par l'ordinateur.

Le pro\-gram\-me
{\idt Cocotte-erroné} ci-dessous contient
de nombreuses erreurs de syntaxe très courantes chez les débutants.

\begin{exemple}[Programme syntaxiquement incorrect]{\hfill

\label{prog:Cocotte-errone} \input Cocotte-errone.tex
\noindent Les erreurs de syntaxe dans ce programme sont les \hbox{suivantes :}
\begin{list}{$\scriptscriptstyle\bullet$}{
\renewcommand{\labelwidth}{3mm}
\renewcommand{\labelsep}{3mm}
\renewcommand{\leftmargin}{10mm}
\renewcommand{\rightmargin}{0mm}
}
\item[1.] le blanc \blanc\ a été oublié entre {\mcl program} et le nom
du programme {\idt Co\-cot\-te$-$er\-ro\-né}$\:$;
\item[2.] le nom du programme {\idt Cocotte$-$erroné} contient un signe moins 
{\idt $-$} et une lettre accentuée {\idt é}$\:$;
\item[3.] la ligne `{\idt {\mcl uses}\ \ldots\ $;$}' a été oubliée$\:$;
\item[4.] l'accolade fermante {\idt \}} a été oubliée à la fin du premier
commentaire$\:$;
\item[5.] l'instruction {\idt pgtg} n'est pas une commande du \robotErr$\:$;
\item[6.] un point-virgule manque après la dernière commande {\idt av}  de la  ligne
commençant par le commentaire {\cmt \{ Dessous \}}$\:$;
\item[7.] le point final a été oublié.
\end{list}
\noindent $\Box$ }\end{exemple}


Un exemple compliqué de programme \pascal\ syntaxiquement correct
est donné ci-\hbox{dessous :}
\begin{exemple}[Le stégosaure]{\rm
Le programme \pascal\ de pilotage du \robot\ pour
dessiner le sté\-go\-sau\-re re\-pré\-sen\-té  à la fi\-gu\-re
\ref{fig:stegosaure} (page \pageref{fig:stegosaure}) est le \hbox{suivant :}
\label{prog:Stegosaure} \input Stegosaure.tex
\noindent $\Box$ }\end{exemple}

\section{Commentaires} 
%\label{sec:4:comment} 

Un programme \pascal\ peut comporter des {\em commentaires}\index{Commentaire} qui sont
des textes français quelconques écrits entre accolades {\cmt \{\ \ldots\ \}}
ou entre pa\-ren\-thè\-ses-étoi\-les {\cmt ($\ast$\ \ldots\ $\ast$)}.
L'or\-di\-na\-teur ne tient pas compte de ces commentaires qui sont uniquement
destinés à expliquer l'usage et le principe de fonctionnement du
programme aux lecteurs de ce programme. Si l'on enlève tous les
commentaires d'un programme, il devient illisible pour un humain mais
rien n'est changé pour l'or\-di\-na\-teur.

\begin{exemple}[Le stégosaure \underline{illisible}]{\rm
En enlevant tous les com\-men\-tai\-res du programme {\idt Stegosaure}
de la page \pageref{prog:Stegosaure}, on obtient un programme \pascal\
totalement illisible mais qui dessine parfaitement bien le 
sté\-go\-sau\-re re\-pré\-sen\-té  à la fi\-gu\-re
\ref{fig:stegosaure} (page \hbox{\pageref{fig:stegosaure}) :}
{\small
{
\parindent \parindentprog
\medskip

\carrenoir
\vspace{\neglskip}\mbox{\blc  \idt  \mcl program\idt \ S$;$
}\mbox{\blc \blk\idt \mcl uses\idt \ \ldots\ $;$
}\par 
\vspace{\neglskip}\mbox{\blc  \idt  \mcl begin\idt \ lc$;$pqtd$;$pqtd$;$av$;$av$;$av$;$av$;$av$;$av$;$av$;$pqtg$;$av$;$av$;$av$;$av$;$av$;$av$;$av$;$av$;$}\par 
\vspace{\neglskip}\mbox{\blc  \idt  av$;$av$;$av$;$av$;$av$;$pqtg$;$bc$;$av$;$pqtd$;$vg$;$pqtg$;$vg$;$pqtd$;$pqtd$;$vg$;$pqtd$;$pqtd$;$vg$;$}\par 
\vspace{\neglskip}\mbox{\blc  \idt  pqtg$;$vg$;$pqtd$;$vd$;$vg$;$av$;$vd$;$av$;$av$;$av$;$pqtd$;$av$;$pqtd$;$av$;$av$;$av$;$av$;$pqtd$;$vd$;$}\par 
\vspace{\neglskip}\mbox{\blc  \idt  av$;$pqtg$;$vd$;$av$;$av$;$av$;$pqtd$;$av$;$pqtd$;$av$;$av$;$pqtd$;$lc$;$av$;$bc$;$vd$;$av$;$av$;$pqtd$;$}\par 
\vspace{\neglskip}\mbox{\blc  \idt  av$;$pqtd$;$av$;$pqtd$;$lc$;$av$;$bc$;$vd$;$pqtg$;$av$;$pqtd$;$av$;$pqtd$;$av$;$av$;$pqtd$;$av$;$av$;$}\par 
\vspace{\neglskip}\mbox{\blc  \idt  pqtd$;$vd$;$pqtg$;$av$;$av$;$vd$;$av$;$pqtd$;$av$;$av$;$pqtd$;$av$;$av$;$pqtd$;$av$;$vd$;$pqtg$;$av$;$av$;$}\par 
\vspace{\neglskip}\mbox{\blc  \idt  pqtd$;$av$;$av$;$pqtd$;$av$;$av$;$pqtd$;$av$;$av$;$pqtd$;$av$;$av$;$pqtg$;$vd$;$av$;$pqtd$;$av$;$av$;$}\par 
\vspace{\neglskip}\mbox{\blc  \idt  pqtd$;$av$;$av$;$pqtd$;$av$;$vd$;$av$;$av$;$pqtg$;$vd$;$pqtd$;$av$;$av$;$pqtd$;$av$;$av$;$pqtd$;$av$;$}\par 
\vspace{\neglskip}\mbox{\blc  \idt  pqtd$;$av$;$pqtg$;$vd$;$vg$;$pqtd$;$pqtd$;$av$;$av$;$pqtd$;$av$;$pqtd$;$av$;$pqtd$;$vg$;$vd$;$vg$;$}\par 
\vspace{\neglskip}\mbox{\blc  \idt  pqtd$;$pqtd$;$av$;$av$;$av$;$pqtd$;$av$;$av$;$pqtd$;$av$;$vd$;$vg$;$av$;$vd$;$pqtd$;$av$;$av$;$av$;$pqtd$;$}\par 
\vspace{\neglskip}\mbox{\blc  \idt  av$;$pqtd$;$av$;$av$;$vd$;$av$;$av$;$lc$;$pqtg$;$av$;$dc$;$pqtd$;$pqtd$;$av$;$pqtg$;$bc$;$vg$;$vd$;$av$;$}\par 
\vspace{\neglskip}\mbox{\blc  \idt  vg$;$pqtg$;$av$;$av$;$pqtd$;$pqtd$;$av$;$av$;$pqtg$;$vg$;$av$;$av$;$av$;$av$;$av$;$av$;$av$;$vg$;$pqtd$;$}\par 
\vspace{\neglskip}\mbox{\blc  \idt  av$;$pqtg$;$av$;$av$;$av$;$pqtd$;$pqtd$;$av$;$av$;$av$;$av$;$av$;$pqtg$;$av$;$av$;$pqtg$;$av$;$av$;$av$;$}\par 
\vspace{\neglskip}\mbox{\blc  \idt  av$;$av$;$av$;$pqtd$;$pqtd$;$av$;$av$;$av$;$av$;$av$;$pqtg$;$av$;$pqtg$;$av$;$av$;$pqtg$;$av$;$pqtd$;$}\par 
\vspace{\neglskip}\mbox{\blc  \idt  pqtd$;$av$;$vg$;$pqtd$;$av$;$av$;$av$;$pqtg$;$av$;$av$;$av$;$pqtd$;$pqtd$;$av$;$av$;$av$;$av$;$av$;$av$;$}\par 
\vspace{\neglskip}\mbox{\blc  \idt  av$;$pqtg$;$av$;$av$;$pqtg$;$av$;$av$;$av$;$av$;$av$;$av$;$av$;$av$;$pqtd$;$pqtd$;$av$;$av$;$av$;$av$;$av$;$}\par 
\vspace{\neglskip}\mbox{\blc  \idt  av$;$av$;$pqtg$;$av$;$pqtg$;$av$;$av$;$av$;$pqtg$;$av$;$pqtd$;$pqtd$;$av$;$av$;$pqtd$;$vg$;$av$;$vd$;$vg$;$}\par 
\vspace{\neglskip}\mbox{\blc  \idt  av$;$pqtd$;$vg$;$pqtd$;$st$;$\mcl end\idt $.$}\par 

\medskip }

}
Cet exemple montre également que la disposition des instructions du
programme est importante. L'espacement entre instructions et les
alignements convenables facilitent grandement la lecture des
programmes.
\noindent $\Box$ }\end{exemple}



\section{Compilation des programmes} 
%\label{sec:4:compil-progs} 
La {\em compilation}\index{Compilation} d'un programme \pascal\ consiste à le traduire en
un {\em programme compilé}\index{Programme compilé} écrit dans le langage de l'or\-di\-na\-teur
appelé {\em langage machine}\index{Langage machine}. La compilation est faite par un programme
que l'on appelle un \index{Compilateur}{\em  compilateur}\@\footnote{\ Nous utiliserons les 
compilateurs Turbo Pascal 5.5 ou 6.0 de Borland sur les or\-di\-na\-teurs compatibles 
IBM PC et le compilateur Turbo Pascal 1.1 de Borland sur les Macintosh.
La façon de compiler les programmes dépend de l'or\-di\-na\-teur
utilisé. Elle est expliquée sur la disquette
d'accompagnement du livre.}.
Bien évidemment le programme \pascal\ et le programme compilé pilotent
le \robot\ de la même façon. Cette manière de procéder nous 
évite d'avoir à apprendre le langage machine. Il est bien difficile
de comprendre ce langage machine car tous les programmes s'écrivent
uniquement avec des 0 et \hbox{des 1$\:$!}

\section{Erreurs syntaxiques} 
%\label{sec:4:err-compil} 

Pour pouvoir compiler un programme \pascal, il ne faut pas qu'il contienne
d'erreurs syntaxiques\index{Erreur de programmation}\index{Erreur syntaxique} violant les règles
d'écriture des programmes. En cas d'erreur le compilateur explique 
brièvement la cause de l'erreur et désigne l'endroit du programme
où l'er\-reur de syntaxe a été trouvée. Il convient alors de
corriger l'erreur et de recommencer la compilation.

\begin{exemple}[Programme syntaxiquement incorrect]{\rm
\hfill Le pro\-gram\-me {\idt Co\-cot\-te$-$\-er\-ro\-né} de la page 
\pageref{prog:Cocotte-errone}
contient à peu près toutes les erreurs syntaxiques imaginables pour
un programme aussi simple.
Quand on compile ce programme {\idt Cocotte$-$erroné}, le compilateur commence par indiquer
`{\mcl program} attendu'. Il lit en effet le mot `{\mcl program}{\idt Cocotte} alors
qu'il s'attend à trouver le mot {\mcl program}.
Ayant ajouté un \blanc\ après le mot
{\mcl program}, une nouvelle tentative de compilation indique 
`{\idt $\scriptstyle ';\scriptstyle '$} attendu', l'erreur étant 
signalée sur le signe moins `{\idt $-$}'.
Diverses causes d'erreurs sont possibles. Nous aurions pu avoir tapé 
moins `{\idt $-$}' à la place du souligné `{\idt \_}' ou 
à la place de  `{\idt $;$}' ou bien par erreur, sans intention de le faire. Le compilateur
qui ne peut pas conna\^\i tre nos véritables intentions, indique simplement la cause d'erreur
la plus pro\-ba\-ble. Ayant remplacé le signe moins `{\idt $-$}' par un
souligné `{\idt \_}', le compilateur signale maintenant une `erreur de syntaxe'
en désignant le `é'. Supprimant l'accent, nous avons corrigé la 
première ligne. Essayant à nouveau de compiler le programme \pascal,
le compilateur signale une erreur par le message `{\mcl begin} attendu' 
en désignant la première commande {\idt  phtd} du \robot. En fait
l'erreur se situe plus tôt, une accolade fermante {\idt \}}
ayant été oubliée à la fin du commentaire expliquant ce que fait 
le programme. Une nouvelle tentative de compilation montre qu'il y a une 
erreur sur la première commande {\idt  phtd} du \robot\ qui est
signalée par le message `identificateur inconnu'. Ayant oublié d'indiquer
que nous utilisons les commandes du \robotErr, le compilateur ignore la
signification de {\idt  phtd}. Après avoir ajouté la ligne
``{\idt {\bf uses}\ \ldots\ $;$}'', nous pouvons tenter une nouvelle compilation.
Cette fois-ci nous obtenons exact\-ement le même message `identificateur 
inconnu' mais le compilateur désigne {\idt  pgtd}. C'est une faute de
frappe que nous corrigeons en remplaçant {\idt  pgtd} par {\idt  phtd}.
Un nouveau message `{\idt ${\scriptstyle '};{\scriptstyle '}$} attendu'
signale qu'un point-virgule manque après l'avant-dernier {\idt  av} 
du programme. L'ayant rajouté, le compilateur indique la dernière
faute `Fin du texte prématurée'. Ayant rajouté le point final nous
obtenons laborieusement le programme {\idt Cocotte} de la page 
\pageref{prog:Cocotte} qui est syntaxiquement correct.
\noindent $\Box$ }\end{exemple}


\section{Exécution des programmes} 
%\label{sec:4:exec-progs} 
L'{\em exécution}\index{Exécution} d'un programme consiste à faire réaliser par
l'or\-di\-na\-teur les actions définies par les instructions du programme\footnote{\ La façon de
demander à l'or\-di\-na\-teur d'exécuter un programme compilé dépend
de l'or\-di\-na\-teur utilisé. Lire les explications avec {\idt LisezMoi}
sur la dis\-quet\-te
d'ac\-com\-pa\-gne\-ment.}. Comme dans notre cas les instructions sont des
commandes du \robot, l'exécution du programme permet de piloter le
\robot. 

Quand le \robot\ dessine trop rapidement pour pouvoir observer le
dé\-rou\-le\-ment du dessin, on peut l'interrompre pour le ralentir ou
l'arrêter à volonté. \index{Exécution (interruption)}
\index{Interruption de l'exécution}Si pendant l'exécution d'un programme de pilotage du \robot, nous
tapons sur la {\em touche d'échappement}\index{Touche d'échappement}
%touche Esc (touche d'echappement)
\begin{picture}(14,12)(0, 2) 
\put(0,0){\framebox(14,12)[lb]{}}
\put(1.5,2){\makebox(0,0)[lb]{\tiny Esc}}
\end{picture}\hspace{0.3ex}
ou sur la touche marquée du
caractère 
%touche §
\begin{picture}(12,12)(0, 2)
\put(0,0){\framebox(12,12)[lb]{}}
\put(1.5,2){\makebox(0,0)[lb]{\tiny \S}}
\end{picture}\hspace{0.3ex}
ou sur celle marquée 
%touche #
\hbox{\begin{picture}(12,12)(0, 2)
\put(0,0){\framebox(12,12)[lb]{}}
\put(1.5,2){\makebox(0,0)[lb]{\tiny \#}}
\end{picture}\hspace{0.3ex},} nous pouvons stopper le dessin du \robot,
exactement comme si nous demandions à l'or\-di\-na\-teur d'exécuter
immédiatement
une commande {\idt st}. Par exemple si nous interrompons l'exécution
du programme {\idt Stegosaure} de la page \pageref{prog:Stegosaure}
au moment où le \robot\ termine le dessin de la patte avant gauche,
nous obtenons l'écran re\-pré\-sen\-té en réduction à la figure
\ref{fig:ecran-stop} de la page \pageref{fig:ecran-stop}.

\begin{figure}[p] %[htbp]
 \includevspacePostScript[7.6cm]{ecran-stop}{0.95}{0.95}
 \caption{Écran de l'or\-di\-na\-teur après un stop}\label{fig:ecran-stop}
\end{figure}

Une fois le \robot\ stoppé, nous \hbox{pouvons :}
\begin{list}{$\scriptscriptstyle\bullet$}{
\renewcommand{\rightmargin}{0mm}
}
\item poursuivre l'exécution du programme normalement en tapant sur la 
touche de retour à la ligne 
\hspace{0.3ex}\framebox[11.mm][l]{$\hookleftarrow$~~~}\hspace{0.3ex},
\item diminuer (en tapant une ou plusieurs fois sur la touche
moins \hbox{%touche -
\begin{picture}(12,12)(0, 2)
\put(0,0){\framebox(12,12)[lb]{}}
\put(1.5,2){\makebox(0,0)[lb]{\tiny $-$}}
\end{picture}\hspace{0.3ex})} ou augmenter (en tapant une ou plusieurs fois sur la 
touche plus \hbox{%touche X
\begin{picture}(12,12)(0, 2)
\put(0,0){\framebox(12,12)[lb]{}}
\put(1.5,2){\makebox(0,0)[lb]{\tiny $+$}}
\end{picture}\hspace{0.3ex}
)} la vitesse du \robot\ (qui peut varier de 0
à 10) puis reprendre l'exécution 
du programme avec cette nouvelle vitesse du \robot\ en tapant sur la 
touche de retour à la ligne 
\hspace{0.3ex}\framebox[11.mm][l]{$\hookleftarrow$~~~}\hspace{0.3ex},
\item arrêter définitivement l'exécution du \robot\ en tapant une 
nouvelle fois sur la  touche d'échappement
\begin{picture}(12,12)(0, 2) %touche d'echappement
\put(0,0){\framebox(12,12)[lb]{}}
\put(1.5,2){\makebox(0,0)[lb]{\tiny esc}}
\end{picture}\hspace{0.3ex} 
ou sur l'une des touches %touche §
\begin{picture}(12,12)(0, 2)
\put(0,0){\framebox(12,12)[lb]{}}
\put(1.5,2){\makebox(0,0)[lb]{\tiny \S}}
\end{picture}\hspace{0.3ex}
ou
%touche #
\begin{picture}(12,12)(0, 2)
\put(0,0){\framebox(12,12)[lb]{}}
\put(1.5,2){\makebox(0,0)[lb]{\tiny \#}}
\end{picture}\hspace{0.3ex}
 (ce qui donne l'écran de la figure
\ref{fig:ecran-arret} page \pageref{fig:ecran-arret}) 
puis en terminant par un retour à la ligne
\hspace{0.3ex}\framebox[11.mm][l]{$\hookleftarrow$~~~}\hspace{0.3ex}.
\end{list}

\begin{figure}[p] %[htbp]
 \includevspacePostScript[7.6cm]{ecran-arret}{0.95}{0.95}
 \caption{Écran de l'or\-di\-na\-teur avant l'arrêt définitif}
 \label{fig:ecran-arret}
\end{figure}

\clearpage
\section{Erreurs logiques} 
%\label{sec:4:err-exec} 

\index{Erreur logique}\index{Erreur de programmation}Après avoir éliminé toutes les erreurs syntaxiques c'est-à-dire 
les erreurs de programmation qui relèvent de la violation des règles de grammaire 
du langage \pascal, nous ne sommes pas encore au bout de nos peines avant de pouvoir
contempler le \hbox{dessin$\:$!} Il reste encore les
{\em erreurs logiques}
que nous découvrirons à l'exécution du programme.

La première erreur possible est que la succession de commandes que notre
programme demande à l'or\-di\-na\-teur de donner au \robot\ conduise celui-ci
en dehors du cadre. C'est le cas du programme suivant :
\input SortirDuCadre.tex
\noindent Ce programme donne le triste résultat 
ci-dessous :
%indiqué à la figure \ref{fig:ecran-sortir} (page \pageref{fig:ecran-sortir}).
\begin{figure}[htbp]
 \includevspacePostScript[7.6cm]{ecran-sortir}{0.95}{0.95}
 \caption{Écran de l'or\-di\-na\-teur quand le robot sort du cadre}
 %\label{fig:ecran-sortir}
\end{figure}

La deuxième erreur possible est que nous nous soyons trompés dans
le programme et que le dessin obtenu ne ressemble pas du tout à ce
que nous avions imaginé. Par exemple, le programme {\idt StegosaureErrone}
ci-dessous est une version logiquement erronée mais syntaxiquement 
correcte du programme {\idt Stegosaure} de la page \pageref{prog:Stegosaure}. 
\input StegosaureErrone.tex
L'auteur de ce programme {\idt StegosaureErrone} a par deux fois confondu sa gauche 
et sa droite. Le 
résultat de l'exécution de ce programme tel
que représenté
ci-dessous
(figure \ref{fig:StegosaureErrone}) %page \pageref{fig:StegosaureErrone} 
éton\-ne\-rait plus d'un pa\-lé\-on\-to\-lo\-gue$\:$!
\begin{figure}[htbp]
 \includehvspacePostScript[105mm,6.8cm]{StegosaureErrone}{0.95}{0.95}
 \caption{Stégosaure plutôt raté}
 \label{fig:StegosaureErrone}
\end{figure}

\section{Comment localiser les erreurs logiques dans le programme} 
%\label{sec:4:trouv-err-exec} 

Pour trouver les erreurs logiques dans le programme il faut s'imaginer
que l'on exécute le programme à la place du \robotErr. Quand le programme 
est très long ceci n'est pas facile.
Le mieux est donc d'arrêter l'exécution du programme un peu avant
l'erreur, puis de sui\-vre l'exécution du programme en mode {\em pas à pas},
\index{Exécution en pas à pas}
ce qui permet de suivre les commandes transmises par l'or\-di\-na\-teur au \robotErr\ pour localiser l'erreur avec précision. Quand l'erreur est
localisée, il faut arrêter l'exécution
du programme pour pouvoir la corriger dans le programme. Une fois
que c'est fait, on peut compiler et tenter une nouvelle
exécution.
{
\newcommand{\touche}[1]{{
\begin{picture}(12,12)(0, 2)
\put(0,0){\framebox(12,12)[lb]{}}
\put(1.5,2){\makebox(0,0)[lb]{\tiny #1}}
\end{picture}\hspace{0.3ex}}}
%
\begin{list}{$\scriptscriptstyle\bullet$}{
\renewcommand{\rightmargin}{0mm}
}
% 1.
\item Pour passer en mode pas à pas nous \hbox{pouvons :}
\end{list}
%
\begin{list}{-}{
\renewcommand{\labelwidth}{3mm}
\renewcommand{\labelsep}{3mm}
\renewcommand{\leftmargin}{15mm}
\renewcommand{\rightmargin}{0mm}
}
%1.1
  \item taper sur la touche d'échappement
    %touche Esc (touche d'echappement)
    \begin{picture}(14,12)(0, 2) 
    \put(0,0){\framebox(14,12)[lb]{}}
    \put(1.5,2){\makebox(0,0)[lb]{\tiny Esc}}
    \end{picture}\hspace{0.3ex},
    %touche §
   \begin{picture}(12,12)(0, 2)
   \put(0,0){\framebox(12,12)[lb]{}}
   \put(1.5,2){\makebox(0,0)[lb]{\tiny \S}}
   \end{picture}\hspace{0.3ex}
   ou
   %touche #
  \begin{picture}(12,12)(0, 2)
  \put(0,0){\framebox(12,12)[lb]{}}
  \put(1.5,2){\makebox(0,0)[lb]{\tiny \#}}
  \end{picture}\hspace{0.3ex}
  puis réduire la vitesse à zéro en tapant plusieurs fois sur la touche
  \touche{$-$} et enfin reprendre l'exécution en mode pas à pas
  en tapant sur la touche de retour à la ligne
  \hspace{0.3ex}\framebox[11.mm][l]{$\hookleftarrow$~~~}\hspace{0.3ex};
% 1.2
  \item ou insérer une instruction {\idt pp} dans le programme \pascal, un
  peu avant l'endroit supposé de l'erreur. Cette instruction {\idt pp} fait passer
  l'or\-di\-na\-teur en mode d'exécution {\underline p}as à {\underline p}as.
\end{list}
%
\begin{list}{$\scriptscriptstyle\bullet$}{
\renewcommand{\rightmargin}{0mm}
}
% 2.
\item En mode pas à pas, l'or\-di\-na\-teur affiche le nom de la commande
du \robot\ qu'il va exécuter puis attend que nous lui donnions
l'ordre de le faire en tapant sur la touche de retour à la ligne
\hspace{0.3ex}\framebox[11.mm][l]{$\hookleftarrow$~~~}\hspace{0.3ex}
(ou en cliquant sur le bouton de la souris). On peut donc suivre
le déroulement de l'exécution du programme instruction après
instruction, ce qui permet en général de trouver assez facilement l'erreur.
% 3.
\item Quand l'erreur est localisée avec précision dans le programme, on
 peut interactivement arrêter l'exécution du
programme en mode pas à pas en tapant sur
 la touche d'é\-chap\-pe\-ment
%touche Esc (touche d'echappement)
\begin{picture}(14,12)(0, 2) 
\put(0,0){\framebox(14,12)[lb]{}}
\put(1.5,2){\makebox(0,0)[lb]{\tiny Esc}}
\end{picture}\hspace{0.3ex},
%touche §
\begin{picture}(12,12)(0, 2)
\put(0,0){\framebox(12,12)[lb]{}}
\put(1.5,2){\makebox(0,0)[lb]{\tiny \S}}
\end{picture}\hspace{0.3ex}
ou
%touche #
\begin{picture}(12,12)(0, 2)
\put(0,0){\framebox(12,12)[lb]{}}
\put(1.5,2){\makebox(0,0)[lb]{\tiny \#}}
\end{picture}\hspace{0.3ex}
puis :
\end{list}
\begin{list}{-}{
\renewcommand{\labelwidth}{3mm}
\renewcommand{\labelsep}{3mm}
\renewcommand{\leftmargin}{15mm}
\renewcommand{\rightmargin}{0mm}
}
% 3.1
   \item reprendre l'exécution normale (en tapant sur la touche
	  \touche{$+$} pour augmenter la vitesse puis  sur la touche de retour à la ligne 
      \hspace{0.3ex}\framebox[11.mm][l]{$\hookleftarrow$~~~}\hspace{0.3ex}),
    \item ou arrêter l'exécution du programme (en tapant à nouveau sur la touche
	  d'é\-chap\-pe\-ment
%touche Esc (touche d'echappement)
\begin{picture}(14,12)(0, 2) 
\put(0,0){\framebox(14,12)[lb]{}}
\put(1.5,2){\makebox(0,0)[lb]{\tiny Esc}}
\end{picture}\hspace{0.3ex},
%touche §
\begin{picture}(12,12)(0, 2)
\put(0,0){\framebox(12,12)[lb]{}}
\put(1.5,2){\makebox(0,0)[lb]{\tiny \S}}
\end{picture}\hspace{0.3ex} ou
%touche #
\begin{picture}(12,12)(0, 2)
\put(0,0){\framebox(12,12)[lb]{}}
\put(1.5,2){\makebox(0,0)[lb]{\tiny \#}}
\end{picture}\hspace{0.3ex}  
  puis sur la touche de retour à la ligne  
  \hbox{\hspace{0.3ex}\framebox[11.mm][l]{$\hookleftarrow$~~~}\hspace{0.3ex}).}
	\end{list}
% 3. (suite)
\begin{list}{$\scriptscriptstyle\bullet$}{
\renewcommand{\rightmargin}{0mm}
}
\item[]  Il est également possible de programmer le retour à une vitesse
   d'\underline{ex}é\-cu\-tion normale en insérant une instruction {\idt ex}
  dans le programme ou de prévoir l'arrêt par {\idt st}.
\end{list}
} %1
Il faut presque toujours plusieurs essais avant que le programme soit 
lo\-gi\-que\-ment correct.
Quand il l'est, on est souvent tenté de l'améliorer. Le plus difficile
est alors de savoir quand \hbox{s'arrêter$\:$!}

\begin{table}[htbp]
\begin{center}
\begin{tabular}{@{}|l|p{11cm}|@{}} \hline
{\idt st}   & \underline{st}oppe momentanément l'exécution du programme 
             de pilotage du \robot. La vitesse peut être diminuée par
 \begin{picture}(12,12)(0, 2) %touche §
  \put(0,0){\framebox(12,12)[lb]{}}
  \put(1.5,2){\makebox(0,0)[lb]{$\scriptscriptstyle -$}}
  \end{picture}\hspace{0.3ex} (minimum 0) ou augmentée par
 \begin{picture}(12,12)(0, 2) %touche §
  \put(0,0){\framebox(12,12)[lb]{}}
  \put(1.5,2){\makebox(0,0)[lb]{$\scriptscriptstyle +$}}
  \end{picture}\hspace{0.3ex} (maximum 10).
			 L'exécution reprend à la vitesse indiquée en tapant sur
			 \hspace{0.3ex}\framebox[11.mm][l]{$\hookleftarrow$~~~}\hspace{0.3ex}
			 ou en cliquant avec la souris. L'exécution  s'arrête
			 définitivement en tapant sur 
			 \begin{picture}(12,12)(0, 2) %touche d'echappement
  \put(0,0){\framebox(12,12)[lb]{}}
  \put(1.5,2){\makebox(0,0)[lb]{\tiny esc}}
  \end{picture}\hspace{0.3ex} ou 
  \begin{picture}(12,12)(0, 2) %touche §
  \put(0,0){\framebox(12,12)[lb]{}}
  \put(1.5,2){\makebox(0,0)[lb]{\tiny \S}}
  \end{picture}\hspace{0.3ex} ou
  \begin{picture}(12,12)(0, 2) %touche #
  \put(0,0){\framebox(12,12)[lb]{}}
  \put(1.5,2){\makebox(0,0)[lb]{\tiny \#}}
  \end{picture}\hspace{0.3ex} puis sur
  \hspace{0.3ex}\framebox[11.mm][l]{$\hookleftarrow$~~~}\hspace{0.3ex}.\\ \hline
{\idt pp}   & exécute la suite du programme en mode \underline{p}as
             à \underline{p}as (vitesse 0). Les commandes du \robot\ s'affichent en bas
		    de l'écran et s'exécutent en tapant sur  
			\hspace{0.3ex}\framebox[11.mm][l]{$\hookleftarrow$~~~}\hspace{0.3ex}.\\ \hline
{\idt ex}   & exécute la suite du programme en mode \underline{ex}écution
             normale (vitesse 10). \\ \hline
\end{tabular}
\end{center}
\label{tab:cmd-vit1}
\caption{Commandes de contrôle du mode de pilotage du \protect\robot}
\end{table}

\noindent \rule{137mm}{0.1mm} %---------------------------------------------
\begin{exercices}{\ \\

\vspace{-\baselineskip}
\noindent \begin{minipage}[t]{88mm}
\tiret Écrire, compiler, mettre au point et exécuter des programmes \pascal\ de dessin du
\robot\ tel qu'il est stylisé ci-contre avec son
crayon baissé et son crayon levé.
\end{minipage}
\hfill
\begin{minipage}[t]{44mm}
\vspace{2.0cm}\includePostScript{Robot+BC+LC}{0.95}{0.95}
\hspace{-4.4cm}  % 4.4cm de large
\end{minipage}

\noindent \begin{minipage}[t]{103mm}
\tiret  Écrire, compiler, mettre au point et exécuter un programme \pascal\ qui
pilote le \robot\ pour lui faire dessiner le chien représenté
ci-contre.
\end{minipage}
\hfill
\begin{minipage}[t]{29mm}
\vspace{2.3cm}\includePostScript{Chien}{0.95}{0.95}
\hspace{-2.9cm}  % 2.9cm de large
\end{minipage}

\noindent \begin{minipage}[t]{103mm}
\tiret  Programmer les dessins de la
pa\-ge \pageref{fig:quadrilateres}.
\end{minipage}

\vspace{1.5mm}

\tiret  Programmer le dessin de vos initiales en lettres majuscules
dans le style de la lettre {\tt H} en relief représentée à la
figure \ref{fig:h-relief} (page \pageref{fig:h-relief}).

\noindent $\Box$}\end{exercices}

\newpage
\begin{corrige}{Nous donnons un programme, les autres se trouvant sur disquette.
\input Chien.tex
\noindent $\Box$}\end{corrige}

%%%%%%%%%%%%%%%%%%%%%%%%%%%%%%%%%%%%%%%%%%%%%%%%%%%%%%%%%%%%%%%%%%%%%%%%%%%
\noindent \framebox[\textwidth]{
\begin{minipage}{\myparboxwidth}
\setlength{\parindent}{2em}
\small \sf
\vspace{\abovedisplayskip}
\section*{Addenda}
Le cycle écriture, mise au point syntaxique, compilation, mise au point
logique et exécution des programmes s'est considérablement
raccourci ces dernières années grâce aux progrès fantastiques
de l'architecture des or\-di\-na\-teurs qu'il convient de comprendre très
schématiquement. 
Un micro-or\-di\-na\-teur \hbox{comprend :}
\begin{list}{$\scriptscriptstyle\bullet$}{
}
\item une unité arithmétique et logique pour faire des calculs et
manipuler des données,
\item une mémoire pour conserver temporairement des données
et des programmes (son contenu est perdu quand on coupe l'électricité),
\item un ou plusieurs disques durs magnétiques 
 pour conserver en permanence des fichiers contenant des données ou des programmes,
\item un ou plusieurs lecteurs de disquettes souples pour
conserver en permanence des fichiers
contenant des données ou des programmes ou les échanger avec d'autres
or\-di\-na\-teurs,
\item un clavier pour entrer des données ou des programmes en mémoire,
\item un écran graphique pour visualiser des textes et des images,
\item une souris pour désigner des points sur l'écran graphique,
\end{list}

\vspace{\abovedisplayskip}
\end{minipage}
}
%%%%%%%%%%%%%%%%%%%%%%%%%%%%%%%%%%%%%%%%%%%%%%%%%%%%%%%%%%%%%%%%%%%%%%%%%%%

%%%%%%%%%%%%%%%%%%%%%%%%%%%%%%%%%%%%%%%%%%%%%%%%%%%%%%%%%%%%%%%%%%%%%%%%%%%
\noindent \framebox[\textwidth]{
\begin{minipage}{\myparboxwidth}
\setlength{\parindent}{2em}
\small \sf
\vspace{\abovedisplayskip}
\begin{list}{$\scriptscriptstyle\bullet$}{
}
\item un bus pour permettre des échanges rapides de données et de
programmes entre ces différents organes.
\end{list}


Quand on met l'or\-di\-na\-teur en marche, il est cablé pour charger dans
sa mémoire et exécuter un programme rangé en permanence sur disque
que l'on appelle le {\em système d'exploitation}\index{Système d'exploitation}. Ce programme a pour
tâche essentielle de gérer tous les organes de l'or\-di\-na\-teur et de
dialoguer avec l'utilisateur grâce à un {\em interpréteur de commandes}\index{Interpréteur de commandes}
qui va attendre que l'utilisateur lui donne des ordres par l'intermédiaire
du clavier et les exécuter immédiatement les uns après les autres
jusqu'à ce que l'utilisateur lui donne l'ordre d'éteindre
l'or\-di\-na\-teur. Parmi les commandes de l'interpréteur, citons par exemple la
suppression d'un fichier rangé sur disque ou le lancement de
l'exécution d'un programme écrit en langage machine (généralement
appelé {\em programme d'application}\index{Programme d'application}) comme par exemple un {\em
éditeur de texte} pour créer un texte et le ranger dans un fichier
ou un {\em compilateur} pour traduire un programme \pascal\ en un programme machine.

Les premières interfaces entre l'homme et la machine étaient purement textuelles. On tape sur
le clavier et on voit le résultat sur une imprimante vite remplacée par un écran avec seulement
des lignes de texte de longueur maximale 80 caractères. Les écrans en noir et blanc furent
remplacés par des écrans couleur.
Puis sont apparus les écrans graphiques constitués de très petits pixels de couleur. Le
fenêtrage permet de donner l'illusion de plusieurs écrans graphiques sur un seul. Puis vint
l'idée d'icônes pour représenter les fichiers et des applications sur lesquels on peut
double cliquer pour lancer une application.

\vspace{\abovedisplayskip}
\end{minipage}
}
%%%%%%%%%%%%%%%%%%%%%%%%%%%%%%%%%%%%%%%%%%%%%%%%%%%%%%%%%%%%%%%%%%%%%%%%%%%

%%%%%%%%%%%%%%%%%%%%%%%%%%%%%%%%%%%%%%%%%%%%%%%%%%%%%%%%%%%%%%%%%%%%%%%%%%%
\noindent \framebox[\textwidth]{
\begin{minipage}{\myparboxwidth}
\setlength{\parindent}{2em}
\small \sf
\vspace{\abovedisplayskip}
Les or\-di\-na\-teurs sont {\em binaires} en ce sens que toutes les données, en
particulier les programmes, sont représentés par des suites de 0 et
de 1.

L'or\-di\-na\-teur fonctionne en transférant une instruction du programme
depuis une case de sa mémoire
dans son unité arithmétique et logique (via le bus) et en 
l'exécutant (ce qui en général modifie le contenu de la mémoire,
du disque, de l'écran, etc.) puis en recommençant ainsi avec l'instruction
rangée dans la case suivante de sa mémoire. Certaines instructions
dites {\em instructions de branchement} permettent d'indiquer que l'instruction suivante
est rangée dans une case mémoire de numéro donné.

\vfill

Pour exécuter un programme écrit en langage machine rangé sur disque,
le système d'exploitation demande à l'unité arithmétique et
logique de charger ce programme en mémoire c'est-à-dire de
recopier le programme dans un endroit de la mémoire
qui était inoccupé puis d'exécuter une instruction de branchement
à la première instruction de ce programme. Le programme doit se terminer
par une instruction de branchement en retour vers le système 
d'exploitation. Les or\-di\-na\-teurs modernes possèdent également la fa\-cul\-té
d'interrompre l'exécution du programme en cours puis de la reprendre
ultérieurement comme si rien ne s'était passé. Entre temps, c'est
le système d'exploitation qui s'exécute pour réagir en fonction de
la cause de l'interruption. Par exemple si l'utilisateur a bougé la souris,
ses nouvelles coordonnées sont rangées en mémoire. Si l'horloge
sonne, ce qui arrive tous les soixantièmes de seconde par exemple, l'heure 
est mise à jour, etc. En particulier l'utilisateur dispose
de boutons ou de touches (%touche Esc (touche d'echappement)
\begin{picture}(14,12)(0, 2) 
\put(0,0){\framebox(14,12)[lb]{}}
\put(1.5,2){\makebox(0,0)[lb]{\tiny Esc}}
\end{picture}\hspace{0.3ex},
%touche §
\begin{picture}(12,12)(0, 2)
\put(0,0){\framebox(12,12)[lb]{}}
\put(1.5,2){\makebox(0,0)[lb]{\tiny \S}}
\end{picture}\nolinebreak\hspace{0.3ex}\nolinebreak,
%touche #
\begin{picture}(12,12)(0, 2)
\put(0,0){\framebox(12,12)[lb]{}}
\put(1.5,2){\makebox(0,0)[lb]{\tiny \#}}
\end{picture}\hspace{0.3ex} 
 pour notre \robot) permettant d'interrompre
l'exécution de son programme s'il constate une erreur.
\vspace{\abovedisplayskip}
\end{minipage}
}
%%%%%%%%%%%%%%%%%%%%%%%%%%%%%%%%%%%%%%%%%%%%%%%%%%%%%%%%%%%%%%%%%%%%%%%%%%%

%%%%%%%%%%%%%%%%%%%%%%%%%%%%%%%%%%%%%%%%%%%%%%%%%%%%%%%%%%%%%%%%%%%%%%%%%%%
\noindent \framebox[\textwidth]{
\begin{minipage}{\myparboxwidth}
\setlength{\parindent}{2em}
\small \sf
\vspace{\abovedisplayskip}
Parmi les applications que peut exécuter le système d'exploitation, 
il y en a une qui nous intéresse particulièrement, c'est le compilateur.

Les compilateurs modernes se chargent de l'édition du texte du
programme, c'est-à-dire de sa création et de ses modifications
par interaction avec l'utilisateur. Le 
programme \pascal\ est affiché sur l'écran et rangé en mémoire 
au fur et à mesure que
l'utilisateur le tape au clavier. Il est prévu qu'il puisse le modifier
en utilisant diverses touches de fonctions et la souris. La possibilité
lui est également offerte de sauvegarder son texte sur le disque dur
dans un fichier dont il doit donner le nom.

Quand l'édition de texte est terminée, le compilateur peut
compiler le programme c'est-à-dire le traduire en un programme équivalent
écrit en langage machine. Le programme compilé peut être rangé
sur disque auquel cas l'utilisateur met fin à l'exécution du compilateur
et retourne au niveau du système d'exploitation pour charger en mémoire
et exécuter son programme compilé comme expliqué précédemment. 
Pendant la phase de débogage du programme, il est plus rapide de demander
au compilateur de ranger le programme compilé en mémoire, le compilateur
se chargeant directement de son exécution puis de reprendre la main quand
cette exécution est terminée ou interrompue.

\vfill

On comprend les progrès considérables qui ont été faits depuis les
premiers or\-di\-na\-teurs pour lesquels il fallait modifier le cablage pour
changer le programme \cite{levy1} :
\begin{list}{$\scriptscriptstyle\bullet$}{
}
\item le programme étant rangé en mémoire (selon l'idée de
Von Neumann\footnote{\ Johann von Neumann, 1903--1957, mathématicien américain.})
et la mémoire reliée à des disques rapides, il est possible de
changer le programme de la machine presque \hbox{instantanément$\:$;}
\item l'usage de langages de haut niveau comme \pascal\ et de compilateurs
performants, permet d'éviter d'avoir à programmer en un langage
machine \hbox{ingrat$\:$;}
\item l'écran alphanumérique permet à l'utilisateur d'éditer son programme
de manière interactive au clavier, ce qui fait gagner un temps 
considérable par rapport aux manipulations de cartes perforées en carton
encore utilisées dans les années \hbox{1970$\:$;}
\item la puissance de calcul, liée à la densité des transitors sur
les circuits (qui double tous les 18 mois ce qui laisse prévoir un milliard
de composants par puce en l'an 2000), permet de gérer des écrans
graphiques directement reliés à la mémoire, ce qui permet des
modifications très rapides et donc l'usage de fenêtres, souris, \hbox{etc$\:$;}
\end{list}
\end{minipage}
}
%%%%%%%%%%%%%%%%%%%%%%%%%%%%%%%%%%%%%%%%%%%%%%%%%%%%%%%%%%%%%%%%%%%%%%%%%%%

%%%%%%%%%%%%%%%%%%%%%%%%%%%%%%%%%%%%%%%%%%%%%%%%%%%%%%%%%%%%%%%%%%%%%%%%%%%
\noindent \framebox[\textwidth]{
\begin{minipage}{\myparboxwidth}
\setlength{\parindent}{2em}
\small \sf
\vspace{\abovedisplayskip}
\begin{list}{$\scriptscriptstyle\bullet$}{
}
\item la taille de la mémoire (quelques millions de caractères) et
des disques (quelques milliards de caractères) permet d'envisager des
programmes et des données de plus en plus grands et donc des applications
de plus en plus sophistiquées.
\end{list}
La taille des plus grands programmes que l'on peut actuellement écrire
se chiffre maintenant en millions de lignes, ce qui est considérable
et ne va pas sans difficultés car on atteint les limites de la
compréhension humaine. Une formation de l'esprit, une 
nouvelle rigueur, plus précise encore que la logique des mathématiciens,
devient la qualité première qui, en plus de l'imagination et du sens de
l'esthétique, fait les bons informaticiens. Un peu de méthode aide
également beaucoup. Toute l'histoire des langages de programmations
est remplie d'inventions permettant de formuler plus clairement la façon
d'organiser les données en mémoire ou les calculs de l'or\-di\-na\-teur. A
chaque fois les capacités humaines, donc la taille limite des plus grands
 programmes, ont été repoussées. 
Une idée très ancienne de structuration des calculs, due à Von Neumann et Turing%
\footnote{\ Alan Turing, 1912--1954, mathématicien anglais.} est celle
de {\em procédure}, on disait autrefois {\em routine}
ou {\em sous-programme}. Nous commençons à l'introduire dans le chapitre suivant.
\vspace{\abovedisplayskip}
\end{minipage}
}
%%%%%%%%%%%%%%%%%%%%%%%%%%%%%%%%%%%%%%%%%%%%%%%%%%%%%%%%%%%%%%%%%%%%%%%%%%%
\newpage
\thispagestyle{empty}